% Данный файл распространяется под лицензией CC BY 4.0. Текст лицензии размещён на https://creativecommons.org/licenses/by/4.0/
% (c) Кафедра системного программирования, 2025

\documentclass[a4paper]{article}
\usepackage[a4paper, top=8mm, bottom=8mm, left=8mm, right=8mm]{geometry}
\usepackage{fontspec}
\usepackage{amsmath}
\usepackage{polyglossia}
\setmainlanguage{russian}
\setotherlanguage{english}
\setmainfont{Liberation Serif} 
\newfontfamily\cyrillicfont{Liberation Serif}
\newfontfamily{\russianfonttt}[Scale=0.7]{DejaVuSansMono}
\usepackage[tiny, compact]{titlesec}
\usepackage{titling}
\setlength{\droptitle}{-1cm}
\pretitle{\begin{center}\begin{bfseries}\Large}
\posttitle{\par\end{bfseries}\end{center}}
\preauthor{\begin{center}\normalsize}
\postauthor{\par\end{center}\vspace{-1.8cm}}

\usepackage{hyperref}
\usepackage{bookmark}
\usepackage{csquotes}
\usepackage{booktabs}

\title{Реализация алгоритмов подсчёта треугольников на основе GraphBLAS API}

\author{Сергеев Дмитрий Сергеевич}

\date{}

\begin{document}

\maketitle

\begin{flushright}
    Группа: 25.Б41-мм

    Кафедра: системного программирования

    Научный руководитель: Григорьев Семён Вячеславович

    Номер семестра практики: 2
\end{flushright}

\section{Введение}
Данная работа выполнена в рамках учебной ознакомительной практики. Основная цель заключается в первичном знакомстве с GraphBLAS API и изучении базовых принципов представления графовых алгоритмов в терминах линейной алгебры. Практика направлена на освоение инструментария и сравнительного анализа производительности существующих библиотечных решений на языках C и Python. 

\section{Постановка задачи}
Целью работы является реализация алгоритмов подсчёта треугольников в графе на основе GraphBLAS API, экспериментальное сравнение собственных реализаций между собой, а также с готовыми реализациями из библиотеки LAGraph.
Для достижения цели требуется решить следующие задачи:
\begin{enumerate}
    \item Выполнить обзор подходов к интеграции GraphBLAS API на языках C и Python.
    \item Реализовать алгоритмы подсчета треугольников с использованием Python-GraphBLAS.
    \item Провести замеры производительности разработанных решений и готовых из LAGraph.
    \item Локализовать причины различий с помощью профилирования.
\end{enumerate}

\section{Обзор существующих решений}
В качестве базы для проведения сравнительного анализа использовались готовые нативные реализации аналогичных алгоритмов из библиотеки LAGraph на языке C. В качестве основного движка как в Python-версии (\texttt{python-graphblas}), так и в сишной версии (\texttt{LAGraph}) выступает библиотека SuiteSparse:GraphBLAS версии 9.4.5.

\section{Описание предлагаемого решения}
Были реализованы три алгоритма подсчёта треугольников с использованием \texttt{python-graphblas}:
\begin{itemize}
    \item \textbf{Алгоритм Burkhard}: основан на формуле $\text{triangles} = \text{sum}(\text{sum}((A^2) \mathbin{.*} A)) / 6$, перемножение матриц $A \times A$ с использованием маски $A$.
    \item \textbf{Алгоритм Sandia}: работает с нижнетреугольной частью матрицы смежности: $\text{triangles} = \text{sum}(\text{sum}((L \times L) \mathbin{.*} L))$, где $L = \text{tril}(A)$.
    \item \textbf{Наивный алгоритм}: вычисляет след матрицы $\text{trace}(A^3) / 6$.
\end{itemize}

\section{Эксперименты}

\subsection*{Тестовый стенд}
Замеры проводились на ноутбуке Huawei MCLF-XX с процессором Intel Core i5-12450H и 16 ГБ оперативной памяти, версия gcc 13.3, python 3.12.3 на Ubuntu 24.04.4. Для стабилизации окружения были приняты следующие меры: отключены фоновые системные сервисы и таймеры, отключён режим Turbo Boost, отключены состояния простоя CPU (C-states), обработка прерываний перенесена на E-cores, замеры привязаны к P-cores через \texttt{taskset -c 0-7}. Замеры производились с питанием от сети и порог заряда был ограничен диапазоном 85-90 процентов. Для каждого алгоритма на каждом графе проводилось 45 запусков, из которых первые 15 считались прогревочными и не учитывались в результатах (для \texttt{python-graphblas}). Перед каждой серией замеров выполнялся сброс файлового кеша ОС.

\subsection*{Наборы данных}
\begin{center}
\begin{tabular}{lrrl}
\toprule
Граф & Вершины & Рёбра\\
\midrule
amazon0302      & 262\,111   & 1\,234\,877 \\
amazon-2008     & 735\,323   & 5\,158\,388  \\
cit-Patents     & 3\,774\,768 & 16\,518\,948  \\
roadNet-CA      & 1\,971\,281 & 2\,766\,607 \\
web-NotreDame   & 325\,729   & 1\,497\,134  \\
web-Stanford    & 281\,903   & 2\,312\,497  \\
\bottomrule
\end{tabular}
\end{center}
Наивный алгоритм был исключён из основных замеров: на графах среднего размера процесс завершался с ошибкой нехватки памяти.

\subsection*{Результаты}
Проводилась статистическая обработка результатов: вычислялись тесты на нормальность, среднее, стандартное отклонение и доверительный интервал с уровнем доверия 95\%.

\begin{table}[h!]
\centering
\caption{Статистические показатели времени выполнения алгоритмов}
\label{tab:benchmarks}
\begin{tabular}{llcccc}
\toprule
\textbf{Graph} & \textbf{Configuration} & \textbf{Mean} & \textbf{Std Dev (\%)} & \textbf{SEM} & \textbf{Confidence Interval 95\%} \\
\midrule
& burkhard-py & 0.1257 & 0.001985 (1.57\%) & 0.000362 & $0.1257 \pm 0.0007$ \\
& sandia-py & 0.03050 & 0.001463 (4.85\%) & 0.000267 & $0.03050 \pm 0.0005$ \\
\cmidrule{2-6}
& burkhard-lagr & 0.1244 & 0.002701 (2.16\%) & 0.000493 & $0.1244 \pm 0.0007$ \\
\textbf{web-NotreDame} & sandia-lagr & 0.0331 & 0.000884 (2.64\%) & 0.000161 & $0.0331 \pm 0.0004$ \\
\midrule
& burkhard-py & 1.1784 & 0.006097 (0.52\%) & 0.001113 & $1.1784 \pm 0.0005$ \\
& sandia-py & 0.2643 & 0.001346 (0.51\%) & 0.000246 & $0.2643 \pm 0.0004$ \\
\cmidrule{2-6}
& burkhard-lagr & 1.1804 & 0.010226 (0.87\%) & 0.001867 & $1.1804 \pm 0.0038$ \\
\textbf{web-Stanford} & sandia-lagr & 0.2607 & 0.002615 (1.00\%) & 0.000477 & $0.2607 \pm 0.0004$ \\
\midrule
& burkhard-py & 0.05783 & 0.000330 (0.57\%) & 0.000060 & $0.05783 \pm 0.00014$ \\
& sandia-py & 0.0464 & 0.000592 (1.30\%) & 0.000108 & $0.0464 \pm 0.0002$ \\
\cmidrule{2-6}
& burkhard-lagr & 0.0564 & 0.000552 (0.99\%) & 0.000101 & $0.0564 \pm 0.0002$ \\
\textbf{roadNet-CA} & sandia-lagr & 0.0543 & 0.000720 (1.32\%) & 0.000131 & $0.0543 \pm 0.0002$ \\
\midrule
& burkhard-py & 4.322 & 0.040947 (0.95\%) & 0.007476 & $4.322 \pm 0.012$ \\
& sandia-py & 0.6794 & 0.002068 (0.30\%) & 0.000378 & $0.6794 \pm 0.0009$ \\
\cmidrule{2-6}
& burkhard-lagr & 4.309 & 0.012695 (0.29\%) & 0.002318 & $4.309 \pm 0.005$ \\
\textbf{cit-Patents} & sandia-lagr & 0.6982 & 0.002949 (0.42\%) & 0.000538 & $0.6982 \pm 0.0012$ \\
\midrule
& burkhard-py & 0.2935 & 0.001290 (0.44\%) & 0.000235 & $0.2935 \pm 0.0004$ \\
& sandia-py & 0.10256 & 0.000299 (0.29\%) & 0.000550 & $0.10256 \pm 0.00013$ \\
\cmidrule{2-6}
& burkhard-lagr & 0.3023 & 0.000890 (0.29\%) & 0.000162 & $0.3023 \pm 0.0003$ \\
\textbf{amazon-2008} & sandia-lagr & 0.1111 & 0.000957 (0.86\%) & 0.000175 & $0.1111 \pm 0.0004$ \\
\midrule
& burkhard-py & 0.08680 & 0.000328 (0.38\%) & 0.000060 & $0.08680 \pm 0.00011$ \\
& sandia-py & 0.02726 & 0.000238 (0.87\%) & 0.000043 & $0.02726 \pm 0.00010$ \\
\cmidrule{2-6}
& burkhard-lagr & 0.0835 & 0.000544 (0.65\%) & 0.000099 & $0.0835 \pm 0.0004$ \\
\textbf{amazon0302} & sandia-lagr & 0.02451 & 0.000876 (3.57\%) & 0.000160 & $0.02451 \pm 0.00013$ \\
\bottomrule
\end{tabular}
\end{table}

1.Алгоритм Sandia быстрее алгоритма Burkhard на всех графах: разница составляет 2-6 раз, доверительные интервалы не пересекаются.

2.Реализация на \texttt{python-graphblas} не уступает реализации \texttt{LAGraph} относительно времени работы. На большинстве графов алгоритмы Burkhard и Sandia показывают статистически неразличимое время, однако на ряде графов \texttt{python-graphblas} отработал быстрее \texttt{LAGraph} (например, на графе \texttt{roadNet-CA}).

В замерах наблюдались выбросы, тест на нормальность не проходился. Тем не менее выбросы присутствуют у обеих реализаций и не меняют относительного положения боксплотов.

\subsection*{Профилирование}

Вклад внутренних функций GraphBLAS рассчитывался вручную. Из логов \texttt{perf} извлекалось абсолютное число собранных сэмплов конкретной исследуемой функции и нормировалось на общее число сэмплов базовой функции, являющейся физической точкой входа в измеряемый цикл вычислений алгоритма.

Профилирование проводилось с частотой сэмплирования 99 Гц для алгоритма Sandia на графах \texttt{roadNet-CA}, \texttt{amazon-2008} и \texttt{web-NotreDame}.

Профилирование показало следующее:
\begin{enumerate}
    \item На графе \texttt{roadNet-CA}: функция \texttt{GB\_AxB\_saxpy3\_flopcount} занимает $\approx 19.08\%$ времени у LAGraph против $\approx 10.41\%$ у \texttt{python-graphblas}; \texttt{GB\_select\_positional\_phase1} -- $\approx 14.96\%$ против $\approx 7.07\%$; \texttt{GB\_msort\_3} -- $\approx 6.59\%$ против $\approx 3.29\%$.
    \item На графе \texttt{amazon-2008} разница между реализациями минимальна (в пределах 1\% на каждую функцию), однако функции, вызываемые этими операциями, также занимают чуть больше времени у LAGraph, что в сумме даёт суммарную разницу.
    \item На графе \texttt{web-NotreDame}: \texttt{GB\_AxB\_saxpy3\_flopcount} -- $\approx 4.83\%$ против $\approx 2.74\%$, \texttt{GB\_msort\_3} -- $\approx 5.22\%$ против $\approx 2.74\%$.
\end{enumerate}

Во всех трёх случаях LAGraph стабильно тратит больше времени на функции: \texttt{GB\_AxB\_saxpy3\_flopcount}, \texttt{GB\_select\_positional\_phase1} и \texttt{GB\_msort\_3}.

\section{Заключение}
В ходе выполнения работы получены следующие результаты:
\begin{itemize}
    \item Реализованы алгоритмы подсчёта треугольников с использованием Python-GraphBLAS API.
    \item Проведено экспериментальное сравнение написанных алгоритмов между собой.
    \item Проведено экспериментальное сравнение реализаций с готовыми алгоритмами из библиотеки LAGraph.
    \item Проведено профилирование.
\end{itemize}

Ознакомиться с построенными боксплотами и флейм графами можно в репозитории: \url{https://github.com/fox56tm/graph_blas_triag_counting}

\end{document}
